% The maximum length of a chess game under the 2023 FIDE Laws is 17,697 plies.
%
% Compile with: latexmk -pdf main_revised_polished_v10.tex
% The bibliography is inline, so no BibTeX pass is required.
\documentclass[11pt]{article}

\usepackage[T1]{fontenc}
\usepackage{lmodern}
\usepackage{amsmath,amssymb,amsthm}
\usepackage{microtype}
\usepackage[margin=3.05cm]{geometry}
\usepackage{needspace}
\usepackage{xcolor}
\usepackage{xurl}
\usepackage[
  colorlinks=true,
  linkcolor=blue!50!black,
  citecolor=blue!50!black,
  urlcolor=blue!50!black
]{hyperref}
\hypersetup{
  pdftitle={The maximum length of a chess game under the 2023 FIDE Laws is 17,697 plies},
  pdfauthor={Junyeop Yim}
}
\urlstyle{same}
\emergencystretch=2em

\newtheorem{theorem}{Theorem}[section]
\newtheorem{lemma}[theorem]{Lemma}
\newtheorem{proposition}[theorem]{Proposition}
\newtheorem{corollary}[theorem]{Corollary}
\theoremstyle{definition}
\newtheorem{definition}[theorem]{Definition}

\newcommand{\B}{\mathtt{B}}
\newcommand{\W}{\mathtt{W}}

\title{The maximum length of a chess game\\
under the 2023 FIDE Laws is 17{,}697 plies}
\author{Junyeop Yim\\
\small Department of Applied Mathematics, Kongju National University\\
\small \texttt{junyeobe0315@smail.kongju.ac.kr}}
\date{10 August 2026}

\begin{document}
\maketitle

\begin{abstract}
The FIDE Laws of Chess effective from 1 January 2023 terminate a game
automatically upon fivefold repetition or after 75 consecutive moves by each
player without a pawn move or a capture. Legal games of $17{,}697$ plies were
previously known, while construction-based scheduling analyses gave an
arithmetic ceiling of $17{,}699$ plies but did not exclude the final two plies.
We close this gap. Partitioning play at pawn moves
and captures yields at most $118$ segments. Each segment has length at most
$150$ plies, and each change in the colour of successive segment endpoints
reduces this bound by one ply. We prove that a game with all $118$ segments
must incur at least three such changes. Hence every game has at most
$150\cdot118-3=17{,}697$ plies, and the known constructions are optimal.
Moreover, in every maximum-length game, all sixteen pawns make six one-rank
moves and promote.
\end{abstract}

\section{Introduction}\label{sec:intro}

How long can a game of chess be under the mandatory termination rules? Before
2014, the answer depended on whether a player chose to claim a draw. Threefold
repetition and the fifty-move rule were claim-based, so cooperating players
could decline to claim and continue indefinitely; classical constructions of
unending play were given by Euwe~\cite{euwe1929} and by Morse and
Hedlund~\cite{morsehedlund1944}. On 1 July 2014, FIDE introduced automatic
draws after fivefold repetition and after 75 moves by each player without a
pawn move or a capture, making the extremal problem finite~\cite{fide2014laws}.
This paper uses the version of the Laws effective from 1 January
2023~\cite{fide2023laws}.

The bookkeeping behind the classical problem was developed by chess
problemists. Bonsdorff, Fabel, and Riihimaa~\cite{bonsdorff1974} summarise the
count of pawn moves and captures that can postpone the fifty-move draw.
Heimo~\cite{heimo2004} observed that replacing 50 by 75 in this count gives
$17{,}697$ plies. After automatic termination was introduced,
Labelle~\cite{labelle2015longest} and Murphy~\cite{murphy2020longest} exhibited
legal games of this length. Murphy's critical-move scheduling analysis produced
an arithmetic ceiling of $17{,}699$ plies and left open whether endpoint parity
forces a deficit of at least three plies in every game with $118$
segments~\cite[Secs.~2.1 and~3]{murphy2020longest}. We prove that it does,
thereby closing the two-ply gap.

Throughout, we disregard time forfeiture and optional endings such as
resignation, agreement, and claimable draws. Play is continued from the
standard initial position until the first mandatory termination condition under
the 2023 Laws applies.

\begin{theorem}\label{thm:main}
Under the FIDE Laws of Chess effective from 1 January 2023, every legal move
sequence from the standard initial position, continued until the first
mandatory termination condition applies, has length at most $17{,}697$ plies.
This bound is attained.
\end{theorem}

Equivalently, a maximum-length game consists of $8{,}848$ complete move pairs
followed by one additional White ply.

The proof of the upper bound reduces to three estimates. Write $L$ for the
game length. A \emph{critical move} is a pawn move or a capture, and critical
moves divide a game into segments. If $K$ is the number of segments and $S$
counts changes in the colour of successive segment endpoints, with the start
of play assigned
Black parity, then
\begin{equation}\label{eq:outline}
  L\le150K-S,\qquad
  K\le118,\qquad
  K=118\Longrightarrow S\ge3.
\end{equation}
The first estimate follows from the seventy-five-move rule and the parity of
segment endpoints. The second couples the number of pawn moves to the number
of pawn captures and then counts the capturable initial pieces. Equality in the
second estimate forces all sixteen pawns to make six moves and at least $29$
captures to occur. These conditions exclude every critical-move block pattern
with at most three blocks and yield the third estimate.

Sections~\ref{sec:setting}--\ref{sec:switches} establish the three estimates in
\eqref{eq:outline}. Section~\ref{sec:maximum} establishes sharpness, completes
the proof, and records the structure forced at equality.

\section{Rules, conventions, and the segment decomposition}
\label{sec:setting}

\subsection{Rule set and chess conventions}

For the purposes of this paper, a \emph{game} is a sequence of legal moves from
the standard initial position, continued until the first mandatory termination
condition in the 2023 FIDE Laws applies. Write $L$ for its length in
\emph{plies}, where one ply is one move by one player; White plays the
odd-numbered plies. The mandatory conditions relevant to the statement are as
follows.

\begin{enumerate}
  \item Checkmate (Article~5.1.1), stalemate (Article~5.2.1), and a dead
    position (Article~5.2.2) terminate the game immediately. A dead position
    is one from which neither player can checkmate by any possible sequence
    of legal moves.
  \item Fivefold repetition (Article~9.6.1), with repeated positions
    interpreted as in Article~9.2, terminates the game.
  \item The seventy-five-move rule (Article~9.6.2) terminates the game after
    75 consecutive moves by each player, equivalently 150 consecutive plies,
    without a pawn move or a capture. If the final ply delivers checkmate,
    checkmate takes precedence.
\end{enumerate}

We disregard time forfeiture, resignation, draws by agreement, and claimable
draws, since none can increase the number of plies. The upper-bound argument
uses Article~9.6.2 throughout and Article~5.2.2 only in
Lemma~\ref{lem:ct30}; the other mandatory termination rules can only shorten a
game. Apart from these rules, the proof uses only alternation of play and the
elementary pawn and capture facts recorded below.

We use the standard coordinates of a chessboard: the files are
$a,b,\ldots,h$ and the ranks are $1,2,\ldots,8$. White pawns begin on rank~2
and move towards rank~8; Black pawns begin on rank~7 and move towards rank~1.
We call ranks~2 and~7 their respective home ranks, and ranks~8 and~1 their
promotion ranks. A pawn advances one rank, or two ranks on its first move from
its home rank, and captures one square diagonally forward. While it remains a
pawn, it changes file only by making such a capture. Promotion upon reaching
the promotion rank is mandatory.

Kings are never captured, so each side has fifteen capturable initial pieces:
eight pawns and seven non-pawn pieces. Throughout the proof, \emph{piece}
includes a pawn. We follow each of the $32$ initial pieces through the game,
retaining its identity after promotion. Thus a promoted pawn remains one of
the eight pieces that began as pawns, and promotion creates no additional
capturable identity.

\subsection{Critical moves and segments}

\begin{definition}[Critical and quiet moves]\label{def:critical}
A \emph{critical move} is a pawn move or a capture. A \emph{quiet move} is a
move that is neither. Thus an en passant capture is critical for both reasons,
promotion is critical because it is a pawn move, and castling is quiet.
\end{definition}

Critical moves are precisely the moves that reset the counter in
Article~9.6.2.

\begin{definition}[Segments and endpoints]\label{def:segments}
Let
\[
  t_1<t_2<\cdots<t_m
\]
be the plies on which the critical moves of a game occur. Set $T=1$ if at
least one quiet ply follows the last critical move, or if the game contains no
critical move, and set $T=0$ otherwise. Define
\[
  e_0=0,\qquad e_i=t_i\quad(1\le i\le m),
\]
and, when $T=1$, define $e_{m+1}=L$. Put $K=m+T$. For
$1\le i\le K$, the interval of plies
\[
  (e_{i-1},e_i]\cap\mathbb Z
  =\{e_{i-1}+1,\ldots,e_i\}
\]
is \emph{segment $i$}. When $T=1$, the final segment is called the
\emph{terminal quiet segment}.
\end{definition}

For $i\ge1$, let $\operatorname{col}(e_i)$ be the colour of the player who
moves on ply $e_i$: White when $e_i$ is odd and Black when it is even. We assign
the virtual endpoint $e_0=0$ the colour Black. This convention treats the start of the
game as having the parity of a preceding Black ply.

\begin{definition}[Switch count]\label{def:switch}
The \emph{switch count} is
\[
  S=\#\bigl\{i\in\{1,\ldots,K\}:
    \operatorname{col}(e_i)\ne\operatorname{col}(e_{i-1})\bigr\}.
\]
\end{definition}

For example, the following endpoint parities give three segments and one
switch:
\[
  0_{\B}\xrightarrow{150}150_{\B}
  \xrightarrow{149}299_{\W}
  \xrightarrow{150}449_{\W}.
\]
The total length is therefore $150\cdot3-1=449$.

\begin{lemma}[Segment decomposition]\label{lem:decomposition}
Every game satisfies
\begin{equation}\label{eq:decomposition}
  L=150K-S-\sum_{i=1}^{K}\delta_i,
  \qquad \delta_i\in\mathbb Z_{\ge0}.
\end{equation}
Here $\delta_i$ is the shortfall of segment $i$ from length $150$ when its
endpoints have the same colour, and from length $149$ when they have different
colours. In particular,
\[
  L\le150K-S.
\]
\end{lemma}

\begin{proof}
Fix a segment with endpoints $e_{i-1}$ and $e_i$. Every ply in the segment
except possibly its last is quiet; the terminal quiet segment is entirely
quiet. If the segment had length at least $151$, then its first $150$ plies
would be quiet, and Article~9.6.2 would terminate the game at ply
$e_{i-1}+150$. Hence every segment has length at most $150$.

Because plies alternate colours, $e_i-e_{i-1}$ is even exactly when the two
endpoints have the same colour. A segment whose endpoints have different
colours therefore has odd length and hence length at most $149$. Subtracting
each segment length from the corresponding bound defines
$\delta_i\in\mathbb Z_{\ge0}$. Summing over the segments gives
\eqref{eq:decomposition}.
\end{proof}

Let $P$ be the number of pawn moves, $C$ the number of captures, and
$C_{\mathrm p}$ the number of captures made by pawns. The number of critical
moves is $P+C-C_{\mathrm p}$, and therefore
\begin{equation}\label{eq:K}
  K=(P-C_{\mathrm p})+(C+T).
\end{equation}

\section{The 118-segment bound}\label{sec:k118}

\begin{lemma}[Pawn moves]\label{lem:pawnbasics}
Each pawn makes at most six pawn moves. Consequently $P\le96$. If a pawn
makes six pawn moves, every one of them advances it by exactly one rank, and
the sixth move promotes it.
\end{lemma}

\begin{proof}
The promotion rank is six rank-steps from the home rank, and every pawn move
advances the pawn by at least one rank towards promotion. Equality therefore
requires six one-rank moves, the last of which promotes the pawn.
\end{proof}

\subsection{Pawn moves and pawn captures}

\begin{definition}[Origin pairs]\label{def:originpair}
For each file, the White pawn and Black pawn that begin on that file form its
\emph{origin pair}. An origin pair is \emph{resolved} if at least one of its
members makes a capture while still a pawn, and \emph{unresolved} otherwise.
A capture made on the promotion move counts as a capture by a pawn; a
capture made by the promoted piece on a later move does not. Let $f$ be
the number of resolved origin pairs.
\end{definition}

\begin{lemma}[Unresolved origin pair]\label{lem:cap10}
The two pawns in an unresolved origin pair make at most ten pawn moves in
total.
\end{lemma}

\begin{proof}
While both pawns remain on the board as pawns, neither changes file, so they
remain on their common origin file with the White pawn strictly below the
Black pawn. If they have advanced $a$ and $b$ ranks, respectively, then
\[
  2+a<7-b,
\]
and hence $a+b\le4$. Thus the pair makes at most four pawn moves while both
members remain on the board. In particular, neither pawn can promote before
the other is captured.

If neither pawn is ever captured, their combined total is at most four.
Otherwise, the pawn captured first has made at most four moves, while the
surviving pawn makes at most six moves over its entire lifetime by
Lemma~\ref{lem:pawnbasics}. The combined total is therefore at most
$4+6=10$.
\end{proof}

\begin{lemma}[Pawn-capture lower bound]\label{lem:pawncapturefloor}
$C_{\mathrm p}\ge f$.
\end{lemma}

\begin{proof}
For each resolved origin pair, choose the first capture made by one of its
pawns. This move contributes to $C_{\mathrm p}$. Different origin pairs give
different moving pawns and hence different moves.
\end{proof}

\begin{proposition}[Net pawn-move bound]\label{prop:pminusc}
For every game,
\[
  P-C_{\mathrm p}\le88.
\]
Equality implies
\[
  f=8,\qquad P=96,\qquad C_{\mathrm p}=8.
\]
\end{proposition}

\begin{proof}
For each unresolved origin pair, Lemma~\ref{lem:cap10} gives at most ten pawn
moves. For each resolved origin pair, Lemma~\ref{lem:pawnbasics} gives at
most twelve. Hence
\[
  P\le\min\bigl\{96,10(8-f)+12f\bigr\}
   =\min\{96,80+2f\}.
\]
Together with $C_{\mathrm p}\ge f$, this gives
\[
  P-C_{\mathrm p}\le\min\{96,80+2f\}-f.
\]
For $f\le7$, the right-hand side is at most $80+f\le87$; for $f=8$, it is at
most $96-8=88$. Equality therefore requires $f=8$, after which the bounds
$P\le96$ and $C_{\mathrm p}\ge8$ force $P=96$ and
$C_{\mathrm p}=8$.
\end{proof}

\subsection{Captures and the terminal quiet segment}

\begin{lemma}[Capture--terminal bound]\label{lem:ct30}
$C+T\le30$.
\end{lemma}

\begin{proof}
There are $30$ capturable initial pieces, so $C\le30$. If $C\le29$, then
$C+T\le30$ because $T\le1$. If $C=30$, only the two kings remain after the
last capture. A position containing only the two kings is dead, so
Article~5.2.2 ends the game on that capture and no terminal quiet segment
follows. Thus $T=0$, and again $C+T=30$.
\end{proof}

\Needspace{14\baselineskip}
\subsection{The segment bound and its equality case}

\begin{theorem}[Segment bound]\label{thm:k118}
Every game satisfies $K\le118$. If $K=118$, then
\begin{equation}\label{eq:equalityprofile}
  f=8,\qquad P=96,\qquad C_{\mathrm p}=8,\qquad C+T=30.
\end{equation}
In particular, all sixteen pawns make six one-rank moves, promoting on the
sixth, and $C\ge29$.
\end{theorem}

\begin{proof}
By \eqref{eq:K}, Proposition~\ref{prop:pminusc}, and Lemma~\ref{lem:ct30},
\[
  K=(P-C_{\mathrm p})+(C+T)\le88+30=118.
\]
Equality requires equality in both component bounds.
Proposition~\ref{prop:pminusc} gives $f=8$, $P=96$, and
$C_{\mathrm p}=8$, while Lemma~\ref{lem:ct30} gives $C+T=30$. Since
$T\le1$, we have $C\ge29$. Finally, $P=96$ forces every pawn to attain the
equality case of Lemma~\ref{lem:pawnbasics}.
\end{proof}

\section{The three-switch bound}\label{sec:switches}

Throughout this section, assume $K=118$. By Theorem~\ref{thm:k118}, every
pawn makes six pawn moves and at least $29$ captures occur. Since $T\le1$, the
game has at least $117$ critical moves; in particular, the critical-move
colour sequence is nonempty.

\subsection{Block patterns}

\begin{definition}[Block pattern]\label{def:pattern}
For a nonempty finite sequence with entries in $\{\B,\W\}$, its \emph{block
pattern} is obtained by replacing each maximal constant run by a single
letter; these maximal runs are the \emph{blocks} of the sequence. The
\emph{critical-move colour sequence} is
\[
  \bigl(\operatorname{col}(t_1),\ldots,
        \operatorname{col}(t_m)\bigr),
\]
and the \emph{endpoint colour sequence} is
\[
  \bigl(\operatorname{col}(e_1),\ldots,
        \operatorname{col}(e_K)\bigr).
\]
Their block patterns are called the \emph{critical-move block pattern} and the
\emph{endpoint block pattern}, respectively.
\end{definition}

For example,
\[
  (\B,\B,\W,\W,\B)\longmapsto\B\W\B.
\]
The virtual endpoint $e_0$ is not part of the endpoint colour sequence; it is
used only in the definition of $S$.

\begin{lemma}[Endpoint blocks]\label{lem:blockrelation}
The endpoint colour sequence has at least as many blocks as the critical-move
colour sequence. If the endpoint colour sequence has $r$ blocks, then
\[
  S=
  \begin{cases}
    r-1,&\text{if the first endpoint is Black},\\
    r,&\text{if the first endpoint is White}.
  \end{cases}
\]
Consequently, an endpoint colour sequence with at least four blocks has
$S\ge3$.
\end{lemma}

\begin{proof}
The endpoint colour sequence is the critical-move colour sequence, with one
terminal endpoint appended when $T=1$. Appending an element cannot reduce the
number of blocks. The formula for $S$ counts the $r-1$ changes between
consecutive blocks and, when the first endpoint is White, the additional
change from the virtual Black endpoint $e_0$.
\end{proof}

The possible critical-move block patterns with at most three blocks are
\[
  \B,\quad \W,\quad \B\W,\quad \W\B,\quad
  \B\W\B,\quad \W\B\W.
\]
We exclude them in turn.

\subsection{One and two blocks}

\begin{proposition}\label{prop:oneblock}
The block patterns $\B$ and $\W$ are impossible when $K=118$.
\end{proposition}

\begin{proof}
A pawn move is a critical move of the pawn's side. With only one block, only
one side's eight pawns can move, so $P\le8\cdot6=48$, contrary to $P=96$.
\end{proof}

\begin{proposition}\label{prop:twoblocks}
The block patterns $\B\W$ and $\W\B$ are impossible when $K=118$.
\end{proposition}

\begin{proof}
Write the pattern as $(X,Y)$, so every critical move by $X$ precedes every
critical move by $Y$. Let $C_X$ and $C_Y$ be the respective numbers of
captures. Since $Y$ can capture at most the fifteen capturable pieces of $X$,
we have $C_Y\le15$, and therefore
\[
  C_X=C-C_Y\ge29-15=14.
\]
All these captures occur before the first critical move by $Y$. Among the
captured $Y$-pieces, at most seven did not begin as pawns, so at least seven
$Y$-pawns are captured before making any pawn move. They cannot move after
they are captured. Hence
\[
  P\le96-7\cdot6=54,
\]
contrary to $P=96$.
\end{proof}

\subsection{A home-rank obstruction}

\begin{lemma}[Home-rank lemma]\label{lem:homerank}
Let $X$ and $Y$ be the two opposing sides. Consider a prefix of a legal game
from the standard initial position. Suppose that
\begin{enumerate}
  \item every move made by $Y$ in this prefix is quiet;
  \item no move made by $X$ in this prefix captures a $Y$-pawn.
\end{enumerate}
Then all eight $Y$-pawns remain on their home squares, no $X$-pawn reaches
$Y$'s pawn home rank, and every $X$-pawn makes at most four pawn moves during
the prefix.
\end{lemma}

\begin{proof}
A move by a $Y$-pawn would be critical, so the first condition implies that
no $Y$-pawn has moved. The second condition implies that none has been
captured. Thus every $Y$-pawn remains on its initial square.

An $X$-pawn cannot move onto $Y$'s pawn home rank. A forward move to an
occupied square is illegal, while a diagonal move to such a square would
capture the $Y$-pawn, contrary to the second condition. Between an $X$-pawn's
home rank and $Y$'s pawn home rank there are four ranks on which it may stand.
Every pawn move advances it by at least one rank, so it makes at most four
pawn moves during the prefix.
\end{proof}

\Needspace{12\baselineskip}
\subsection{Three blocks}

The remaining case is ruled out by the home-rank obstruction.

\begin{proposition}\label{prop:threeblocks}
The block patterns $\B\W\B$ and $\W\B\W$ are impossible when $K=118$.
\end{proposition}

\begin{proof}
Write the pattern as $(X,Y,X)$, and let $C_X,C_Y$ be the numbers of captures
made by the two sides. Since $X$ can capture at most fifteen $Y$-pieces,
\[
  C_Y=C-C_X\ge29-15=14.
\]
At most seven of the pieces captured by $Y$ did not begin as pawns. Hence at
least seven of the captured $X$-pieces began as pawns.

Let $\tau$ be the first critical move by $Y$, and fix one such piece $p$,
which began as an $X$-pawn. By Theorem~\ref{thm:k118}, $p$ makes six pawn
moves before it is captured. Any critical move by $X$ after $\tau$ belongs to
the final $X$-block and therefore occurs after every critical move in the
middle $Y$-block, including the capture of $p$. Consequently all six pawn
moves of $p$ occur before $\tau$.

Consider the game prefix ending immediately before $\tau$. Every move made by
$Y$ in this prefix is quiet. Moreover, $X$ cannot have captured a $Y$-pawn:
such a pawn would have been removed before making any critical move and could
not make the six pawn moves required by $P=96$. Lemma~\ref{lem:homerank}
therefore applies to this prefix and implies that $p$ makes at most four pawn
moves before $\tau$, contradicting the six moves just established.
\end{proof}

\begin{theorem}[Three-switch bound]\label{thm:s3}
If $K=118$, then the critical-move colour sequence has at least four blocks
and $S\ge3$.
\end{theorem}

\begin{proof}
Propositions~\ref{prop:oneblock},~\ref{prop:twoblocks}, and~\ref{prop:threeblocks}
exclude all nonempty block patterns with at most three
blocks. Hence the critical-move colour sequence has at least four blocks.
Lemma~\ref{lem:blockrelation} transfers this lower bound to the endpoint
colour sequence and gives $S\ge3$.
\end{proof}

\section{Completion of the proof and the equality case}
\label{sec:maximum}

\Needspace{11\baselineskip}
\begin{proposition}[Sharpness]\label{prop:sharpness}
There exists a legal move sequence of $17{,}697$ plies from the standard
initial position such that no mandatory termination condition applies before
the final ply, and the final ply delivers checkmate.
\end{proposition}

\begin{proof}
Labelle~\cite{labelle2015longest} and Murphy~\cite{murphy2020longest} gave
explicit legal games of this length. For reproducibility, the companion
archive~\cite{repo} contains Murphy's move sequence and two separately
implemented replay verifiers. The verifiers check move legality, fivefold
repetition, and the seventy-five-move rule after every ply, and confirm
checkmate on the final ply. Neither implements a complete decision procedure
for dead positions under Article~5.2.2, but no such procedure is needed for
this witness. From every earlier position, the remaining verified suffix is a
legal sequence ending in checkmate, so that position is not dead; the existence
of the next legal ply also rules out earlier checkmate and stalemate. Hence no
mandatory termination condition applies before the final ply, and the lower
bound follows.
\end{proof}

\begin{proof}[Proof of Theorem~\ref{thm:main}]
Let $K$ be the number of segments. By Theorem~\ref{thm:k118}, $K\le118$. If
$K\le117$, then Lemma~\ref{lem:decomposition} gives
\[
  L\le150\cdot117=17{,}550.
\]
If $K=118$, then Theorem~\ref{thm:s3} gives $S\ge3$, and therefore
\[
  L\le150\cdot118-3=17{,}697.
\]
This proves the universal upper bound. Proposition~\ref{prop:sharpness} shows
that it is attained.
\end{proof}

\begin{corollary}[Structure of maximum-length games]\label{cor:equality}
Every game of length $17{,}697$ has
\[
  K=118,\qquad S=3,\qquad \delta_i=0\quad(1\le i\le118),
\]
and
\[
  f=8,\qquad P=96,\qquad C_{\mathrm p}=8,\qquad C+T=30.
\]
Thus all sixteen pawns make six one-rank moves, promoting on the sixth; at
least $29$ captures occur; and every segment has the maximum length permitted
by its endpoint parity.
The critical-move and endpoint block patterns are both $\B\W\B\W$.
\end{corollary}

\begin{proof}
A game with $K\le117$ has length at most $17{,}550$, so a maximum-length game
has $K=118$. Combining Theorem~\ref{thm:s3} with
\eqref{eq:decomposition} gives
\[
  17{,}697
  =150\cdot118-S-\sum_{i=1}^{118}\delta_i
  \le150\cdot118-3.
\]
Equality holds throughout. Hence $S=3$ and every $\delta_i=0$. The equality
statement in Theorem~\ref{thm:k118} gives the remaining numerical assertions
and the behaviour of the pawns.

The critical-move colour sequence has at least four blocks by
Theorem~\ref{thm:s3}. Each transition between consecutive blocks contributes
a switch between critical endpoints, so $S=3$ permits at most four blocks.
Thus there are exactly four. If the first block were White, the comparison
with the virtual Black endpoint would contribute a fourth switch. Hence the
critical-move block pattern is $\B\W\B\W$. Appending a terminal endpoint
cannot create a new block without increasing $S$, so the endpoint block
pattern is the same.
\end{proof}

\section*{Data and code availability}

The companion archive~\cite{repo} contains Murphy's $17{,}697$-ply move
sequence together with two separately implemented replay verifiers. The
verifiers check move legality, fivefold repetition, and the seventy-five-move
rule after every ply, and confirm checkmate on the final ply. The proof of
Proposition~\ref{prop:sharpness} explains why the verified suffix ending in
checkmate also rules out dead positions for this witness. The archive also
contains ancillary computational cross-checks. None of these computations
enters the upper-bound proof.

The pinned software snapshot supporting the witness verification and ancillary
checks is version~1.1.0 at Git commit
\nolinkurl{2000be4d71d3e5ea72e565792c559a03e3dbc5fe}. The cited Zenodo DOI is
the all-versions record; the commit hash fixes the exact software revision. The
witness
is \nolinkurl{data/longest.pgn}, whose SHA-256 digest is
\nolinkurl{6700b7b70260c9b4448d58c610601cab938dd0e01392bd876d3379630de680bb}.
From a clean checkout, \texttt{make witness} replays the witness and checks the
expected length $17{,}697$ and termination by checkmate. The command
\texttt{make verify} additionally runs the counting checks, test suite,
independent move-generator check, and model cross-checks. The move-generator
and CP-SAT cross-checks degrade to skips when their optional C compiler or
OR-Tools dependency is unavailable. The development version and detailed
verification notes are maintained in the
\href{https://github.com/junyeobe0315/long_chess}{project repository on
GitHub}.

\section*{Acknowledgements}

Tom Murphy VII's construction and question motivated this work. The author
thanks Fran\c{c}ois Labelle for documenting the earlier constructions and
Alexis Langlois-R\'emillard for comments on an earlier draft.

\begin{thebibliography}{9}

\bibitem{bonsdorff1974}
E.~Bonsdorff, K.~Fabel, and O.~Riihimaa,
\emph{Schach und Zahl: Unterhaltsame Schachmathematik},
2nd ed., Walter Rau Verlag, D\"usseldorf, 1971.
Spanish translation: \emph{Ajedrez y matem\'aticas},
Ediciones Mart\'inez Roca, Barcelona, 1974.

\bibitem{euwe1929}
M.~Euwe,
Mengentheoretische Betrachtungen \"uber das Schachspiel,
\emph{Proceedings of the Koninklijke Akademie van Wetenschappen te Amsterdam}
\textbf{32} (1929), 633--642.

\bibitem{fide2014laws}
FIDE,
\emph{Laws of Chess: For competitions starting from 1 July 2014 till
30 June 2017},
FIDE Handbook E/01, 2014,
\url{https://handbook.fide.com/chapter/E012014}.

\bibitem{fide2023laws}
FIDE,
\emph{FIDE Laws of Chess: Taking effect from 1 January 2023},
FIDE Handbook E/01, 2023,
\url{https://handbook.fide.com/chapter/e012023}.

\bibitem{heimo2004}
O.~Heimo,
Pisin shakkipeli / The longest chess game,
\emph{Suomen Teht\"av\"aniekat} \textbf{58} (2004), no.~4--5, 159--162,
\url{https://www.tehtavaniekat.fi/wp-content/uploads/vanhatlehdet/ST_2004_4-5.pdf}.

\bibitem{labelle2015longest}
F.~Labelle,
The longest possible chess game, and bounds on the number of possible chess
games, 2015, updated 2020,
\url{https://wismuth.com/chess/longest-game.html}.

\bibitem{morsehedlund1944}
M.~Morse and G.~A.~Hedlund,
Unending chess, symbolic dynamics and a problem in semigroups,
\emph{Duke Mathematical Journal} \textbf{11} (1944), no.~1, 1--7.

\bibitem{murphy2020longest}
T.~Murphy VII,
Is this the longest chess game?,
in \emph{A Record of the Proceedings of SIGBOVIK 2020}, 2020, pp.~20--31,
\url{https://tom7.org/chess/longest.pdf}.

\bibitem{repo}
J.~Yim,
\emph{long\_chess: companion repository for ``The maximum length of a chess
game under the FIDE Laws is 17,697 plies''},
software version~1.1.0, Git commit
\nolinkurl{2000be4d71d3e5ea72e565792c559a03e3dbc5fe},
Zenodo, 2026, all-versions DOI:
\url{https://doi.org/10.5281/zenodo.21828025}.
Exact software snapshot:
\url{https://github.com/junyeobe0315/long_chess/tree/2000be4d71d3e5ea72e565792c559a03e3dbc5fe}.

\end{thebibliography}

\end{document}
